\documentclass[11pt,a4paper]{article}


\usepackage[margin=2cm]{geometry}
\usepackage{biblatex}
\usepackage{amsmath,amssymb,amsthm}
% \usepackage{amsthm}
% \usepackage[margin=2cm]{geometry}
\usepackage{dsfont}
\usepackage{tikz}
\usetikzlibrary{cd}
\usetikzlibrary{babel,calc}
\usetikzlibrary{positioning,arrows.meta,fit,backgrounds}
% \usepackage{tikz-cd}
% \usepackage{titlesec}
\usepackage{enumitem}
\usepackage{hyperref}
\usepackage{cleveref}
% \usepackage{parskip}
\usepackage{ebproof}
% \usepackage{quiver}
\usepackage{stackengine,scalerel}
% \usepackage{xparse}
% \usepackage{physics}
% \usepackage{physics2}
% \usephysicsmodule{ab,ab.braket}
% \usepackage{tikzit}
% \input{smc.tikzstyles}
\usepackage{xcolor}
\usepackage{graphicx}
\usepackage{array}    % m{} columns + \arraybackslash (combined car-model table)
\usepackage{booktabs} % \toprule/\midrule/\bottomrule

\usepackage{mathtools}
\usepackage{stmaryrd}
\usepackage{subcaption}

\newtheorem*{fact}{Fact}

\crefname{prop}{prop.}{props.}
\Crefname{prop}{Proposition}{Propositions} % Capitalized version for start of sentences

\crefname{remark}{remark}{remarks}
\Crefname{remark}{Remark}{Remarks} % Capitalized version for start of sentences

\crefname{appendix}{appendix}{appendices}
\Crefname{appendix}{Appendix}{Appendices} % Capitalized version for start of sentences

% \crefname{property}{property}{properties}
% \Crefname{property}{Property}{Properties} % Capitalized version for start of sentences

\crefformat{property}{#2#1#3}
\Crefformat{property}{#2#1#3}

\crefname{def}{definition}{definitions.}
\Crefname{def}{Definition}{Definitions} % Capitalized version for start of sentences

\crefformat{section}{#2\S#1#3}
\Crefformat{section}{#2\S#1#3}

\usepackage{listings}

\usepackage{tcolorbox}
\tcbuselibrary{listings} % Import listings support

\definecolor{codegreen}{rgb}{0,0.6,0}
\definecolor{codegray}{rgb}{0.5,0.5,0.5}
\definecolor{codepurple}{rgb}{0.58,0,0.82}
\definecolor{backcolour}{rgb}{0.95,0.95,0.92}

\lstdefinestyle{mystyle}{
    backgroundcolor=\color{backcolour},
    commentstyle=\color{codegreen}\textit,
    keywordstyle=\color{blue},
    numberstyle=\tiny\color{codegray},
    stringstyle=\color{codepurple},
    basicstyle=\ttfamily\footnotesize,
    breakatwhitespace=false,
    breaklines=false,
    captionpos=b,
    keepspaces=true,
    numbers=left,
    numbersep=5pt,
    showspaces=false,
    showstringspaces=false,
    showtabs=false,
    tabsize=2
}

\lstset{style=mystyle}

\tcbset{
    mystyle/.style={
        listing only,
        breakable, % Allows box to span pages
        colback=CodeBack,
        colframe=CodeFrame,
        boxrule=0.5pt,
        sharp corners,
        listing options={
            backgroundcolor=\color{backcolour},   
            commentstyle=\color{codegreen}\textit,
            keywordstyle=\color{blue},
            numberstyle=\tiny\color{codegray},
            stringstyle=\color{codepurple},
            basicstyle=\ttfamily\footnotesize,
            breakatwhitespace=false,         
            breaklines=true,                 
            captionpos=b,                    
            keepspaces=true,                 
            numbers=left,                    
            numbersep=5pt,                  
            showspaces=false,                
            showstringspaces=false,
            showtabs=false,                  
            tabsize=2,
        }
    }
}

\newtcblisting{codelisting}[3][]{
    listing only,
    % breakable, % Allows box to span pages
    colback=backcolour,
    % colframe=CodeFrame,
    boxrule=0pt,
    boxsep=0pt,
    top=1mm,
    bottom=1mm,
    sharp corners,
    notitle,
    listing options={
        language=#2,
        morekeywords= {exo, endo},
        commentstyle=\color{codegreen}\textit,
        keywordstyle=\color{blue},
        numberstyle=\tiny\color{codegray},
        stringstyle=\color{codepurple},
        basicstyle=\ttfamily\tiny,
        breakatwhitespace=false,         
        breaklines=true,                 
        captionpos=b,                    
        keepspaces=true,                 
        numbers=left,                    
        numbersep=5pt,                  
        showspaces=false,                
        showstringspaces=false,
        showtabs=false,                  
        tabsize=2,
        aboveskip=0pt,  % Removes the space above the code
        belowskip=0pt,  % Removes the space below the code
        mathescape=#3,
        #1               % optional extra listings options (e.g. escapeinside for \cnew)
    }
}

% Highlight code that is new relative to the previous rung's program.
% Use inside a codelisting that passes escapeinside, e.g.
%   \begin{codelisting}[escapeinside={(*}{*)}]{haskell}{true} ... (*\cnew{new}*) ...
\definecolor{codenew}{RGB}{255,228,130}
\newcommand{\cnew}[1]{{\setlength{\fboxsep}{0.6pt}\colorbox{codenew}{\ttfamily #1}}}
% Coloured (not highlighted) font, for marking new syntax in the grammar figure.
\newcommand{\cnewm}[1]{\textcolor{blue!65!black}{#1}}

% Theo's Macros
\newcommand{\naturals}{\mathbb{N}}
\newcommand{\integers}{\mathbb{Z}}
\newcommand{\booleans}{\mathbb{B}}
\newcommand{\indicator}{\mathds{1}}
\newcommand{\va}{\mathbf{a}}
\newcommand{\vb}{\mathbf{b}}
\newcommand{\vx}{\mathbf{x}}
\newcommand{\vy}{\mathbf{y}}
\newcommand{\vw}{\mathbf{w}}
\newcommand{\vz}{\mathbf{z}}
\newcommand{\vsigma}{\mathbf{\sigma}}
\newcommand{\reals}{\mathbb{R}}
\newcommand{\complex}{\mathbb{C}}
\newcommand{\field}{\mathbb{F}}
\newcommand{\sign}{\mathrm{sign}}
\DeclareMathOperator*{\argmin}{arg\,min}
\newcommand{\true}{\mathsf{True}}
\newcommand{\false}{\mathsf{False}}
\newcommand{\brackets}[1]{\{#1\}}

% \newcommand{\bicatb}{\mathcal{B}}

\newcommand{\catc}{\mathbf{C}}
\newcommand{\catcop}{\catc^{op}}
\newcommand{\catd}{\mathbf{D}}
\newcommand{\cate}{\mathcal{E}}
\newcommand{\eqdef}{\triangleq}
\newcommand{\id}{\mathbf{id}}
% \newcommand{\idbicat}{\mathbf{Id}}
\newcommand{\Id}{\mathbf{Id}}
\newcommand{\righttop}[1]{\xrightarrow{#1}}
\newcommand{\lefttop}[1]{\xleftarrow{#1}}
\newcommand{\Ob}{\mathbf{Ob} \ }
\newcommand{\Hom}{\mathrm{Hom}}
% \newcommand{\yo}{\textmc{よ}}
\newcommand{\yo}{\textbf{y}}

\newcommand{\terminalmorph}{\langle\rangle}
\newcommand{\prodmorph}[2]{\langle#1, #2\rangle}

\newcommand{\ul}[1]{\underline{#1}}
\newcommand{\ulyo}{\ul{\yo}}
\newcommand{\yozero}{\yo(0)}
\newcommand{\yoone}{\yo(1)}
\newcommand{\yon}{\yo(n)}
\newcommand{\yom}{\yo({m})}
\newcommand{\yop}{\yo(p)}
% \newcommand{\op}{\mathbf{op}}
\newcommand{\Ops}{\mathcal{O}}
% \newcommand{\typing}[3]{#1 \ | \ #2 \vdash #3}
% \newcommand{\typingfilled}[1]{\Gamma \ | \ \Delta \vdash #1}
% \newcommand\eye{\ensurestackMath{\stackinset{c}{}{c}{-.33pt}%
%   {\bullet}{\bigcirc}}}
% \newcommand\smallereye{\scalerel*{\eye}{x}}
% \newcommand{\eyemid}{ \ \eye \ }

% Vectors
\newcommand{\veca}{\Vec{a}}
\newcommand{\vecb}{\Vec{b}}
\newcommand{\vecc}{\Vec{c}}
\newcommand{\vecd}{\Vec{d}}
\newcommand{\vecx}{\Vec{x}}
\newcommand{\vecy}{\Vec{y}}
\newcommand{\vecz}{\Vec{z}}

\newcommand{\alphaequiv}{\simeq_\alpha}

\newcommand{\LHS}{\mathsf{LHS}}
\newcommand{\RHS}{\mathsf{RHS}}

% Chapter 2 algebraic theories
\newcommand{\Mon}{\mathsf{Mon}}

\newcommand{\Tree}{\mathsf{Tree}}

\newcommand{\Kl}[1]{\mathsf{Kl}({#1})}

\newcommand{\Name}{\mathcal{N}}
\newcommand{\Var}{\mathsf{Var}}
\newcommand{\IPVar}{\mathsf{IPVar}}
\newcommand{\GroundValue}{\mathsf{GroundValue}}
\newcommand{\Value}{\mathsf{Value}}
\newcommand{\Expr}{\mathsf{Expr}}

\newcommand{\bind}{\mathbin{>\!\!\!>\mkern-6.7mu=}}
\newcommand{\fin}{\mathsf{fin}}
\newcommand{\MVal}{\mathsf{MVal}}
\newcommand{\Intervention}{\mathsf{Intv}}
\newcommand{\MultiWorld}{\mathsf{MultiWorldCounterfactual}}
\newcommand{\TwinWorld}{\mathsf{TwinWorldCounterfactual}}
\newcommand{\supp}{\mathsf{supp}}

% ------------- Formal Def ----------------
\newcommand{\IdealisedChiRho}{\mathsf{ChiRho}}
\newcommand{\Prob}{\mathsf{Prob}}
\newcommand{\Dist}{\mathsf{Dist}}
\newcommand{\Paths}{\mathsf{Paths}}
\newcommand{\UPaths}{\mathsf{UPaths}}
\newcommand{\posreals}{\reals_{\geq 0}}

% Syntax
% \newcommand{\opif}[3]{\mathbf{if} \ #1 \ #2 \ #3}
% \newcommand{\lift}[1]{\overline{#1}}
\newcommand{\opprod}{\mathsf{prod}}
\newcommand{\opscore}{\mathsf{score}}
\newcommand{\opflip}{\mathsf{flip}}
\newcommand{\dbernoulli}{\mathbf{bernoulli}}
\newcommand{\dnormal}{\mathbf{normal}}
\newcommand{\constr}{\mathbf{constr}}
\newcommand{\intvnone}{\mathsf{None}}
\newcommand{\intvsome}{\mathsf{Some}}
\newcommand{\opskip}{\mathsf{skip}}
\newcommand{\opsample}{\mathsf{sample}}
\newcommand{\opintervene}{\mathsf{intervene}}
\newcommand{\opbroadcast}{\mathsf{broadcast}}
\newcommand{\optrack}{\mathsf{track}}
\newcommand{\opdistribute}{\mathsf{distribute}}
\newcommand{\opif}{\mathbf{if}}
\newcommand{\ifelse}[3]{\opif \ #1 \ \mathbf{then} \ #2 \ \mathbf{else} \ #3}
\newcommand{\opthunk}{\mathsf{thunk}}
\newcommand{\opforce}{\mathsf{force}}
\newcommand{\returnSymb}{\mathbf{return}}
\newcommand{\return}[1]{\returnSymb \ #1}
\newcommand{\returnParam}[2]{\returnSymb_{#1} \ #2}
\newcommand{\app}{\odot}
\newcommand{\pureval}{\mathbf{pureVal}}
\newcommand{\mem}{\mathsf{mem}}
\newcommand{\letbind}[3]{\mathbf{let} \ #1 = #2 \ \mathbf{in} \ #3}
\newcommand{\opdo}{\mathsf{do}}
\newcommand{\rhide}{\mathsf{hide}}
\newcommand{\rwk}{\mathsf{wk}}
\newcommand{\rcontr}{\mathsf{contr}}
\newcommand{\rexch}{\mathsf{exch}}
\newcommand{\rrename}{\mathsf{rename}}
\newcommand{\Renamings}{\mathsf{Renamings}}
\newcommand{\Deltahat}{\widehat{\Delta}}
\newcommand{\projf}{\pi_{\mathsf{f}}}
\newcommand{\projcf}{\pi_{\mathsf{cf}}}
\newcommand{\coerce}[3]{\mathsf{coerce}_{#2}^{#1}[#3]}
\newcommand{\rv}{\underline{r}}
\newcommand{\restr}{\rvert}
\newcommand{\subst}[2]{[#1/#2]}
\newcommand{\ket}[1]{| #1 \rangle}

% Types
\newcommand{\tnat}{\mathsf{Nat}}
\newcommand{\tbool}{\mathsf{Bool}}
\newcommand{\tint}{\mathsf{Int}}
\newcommand{\tposreal}{\mathsf{PosReal}}
\newcommand{\treal}{\mathsf{Real}}
\newcommand{\tunit}{\mathsf{Unit}}
\newcommand{\tprod}{\mathsf{Prod}}
\newcommand{\tP}[2]{P \ #1 \ #2}
\newcommand{\unitgrade}{[]}
\newcommand{\rarrow}{\Rightarrow}

% Type system
\newcommand{\dtyping}[3]{#1 \vdash^d #2: #3}
\newcommand{\ptyping}[4]{#1 \mid #2 \vdash^p #3: #4}
\newcommand{\rtyping}[4]{#1 \vdash #2: #3 \rarrow #4}

\newcommand{\ctxequiv}{\simeq_{\mathsf{ctx}}}
\newcommand{\Values}{\mathsf{Values}}

% DenSem
% \newcommand{\sem}[1]{\llbracket #1 \rrbracket}
\newcommand{\subfin}{\subseteq_{\fin}}
\newcommand{\inj}{\mathbf{inj}}
\newcommand{\detsem}[1]{[#1]}
\newcommand{\List}{\mathrm{List}}
\newcommand{\sIntervention}{\mathrm{Intv}}
\newcommand{\sP}[2]{\mathbf{P} \ #1 \ #2}
\newcommand{\sPP}{\mathbf{P}}
\newcommand{\sVal}{\mathbf{Val}}
\newcommand{\gradecatsymb}{\mathbb{G}}
\newcommand{\gradecat}[1]{\gradecatsymb_{#1}}
\newcommand{\sdist}{\mathrm{dist}}
\newcommand{\sintervene}{\mathrm{intervene}}

\DeclareMathOperator*{\dom}{dom}

% New language
\newcommand{\sCaus}{\mathrm{Caus}}
\newcommand{\Caus}[2]{\sCaus \ #1 \ #2}
% \newcommand{\Maybe}{\mathsf{Maybe}}
\newcommand{\opnew}{\mathsf{new}}
\newcommand{\Caustic}{\mathsf{\lambda_C}}
\newcommand{\cdtyping}[3]{#1 \vdash^d_{\Caustic} #2: #3}
\newcommand{\crtyping}[4]{#1 \vdash^r_{\Caustic} #2: #3 \rarrow #4}
\newcommand{\cptyping}[4]{#1 \vdash^{p}_{\Caustic} #2 :#3 \& #4}
\newcommand{\cptypingreordered}[4]{\cptyping{#2}{#3}{#4}{#1}}

% Compilation and Logical Relation
\DeclareMathOperator*{\len}{dom}

\newcommand{\Maybe}[1]{\mathrm{Maybe}({#1})}
\newcommand{\sem}[1]{\llbracket{#1}\rrbracket}
\newcommand{\Just}[1]{\mathrm{Just}({#1})}
\newcommand{\Nothing}{\mathrm{Nothing}}

% Dario's Macros
\usepackage{bm}
\renewcommand{\vec}{\bm}

\newcommand{\todo}[1]{\textcolor{blue}{[TODO: #1]}}
\newcommand{\prompt}[1]{\textcolor{purple}{[PROMPT: #1]}}

\newcommand{\cat}[1]{\mathbb{#1}}
\newcommand{\catname}[1]{\mathbf{#1}}

\newcommand{\set}{\mathbf{set}}
\newcommand{\Set}{\catname{Set}}

\newcommand{\concat}{+\!\!+}

\newcommand{\defeq}{\coloneq}


\newcommand{\C}{\cat C}
\newcommand{\G}{\mathcal G}

\newcommand{\op}{\mathrm{op}}


\newcommand{\thought}[3]{\marginpar{\tiny\itshape\raggedright\textcolor{#2}{\textbf{#1:} #3}}\ignorespaces} \marginparwidth=1.4cm \marginparsep=0cm
\newcommand{\jeremy}[1]{\thought{Jeremy}{green!20!black}{#1}}

\newcommand{\IntvInstr}{\mathsf{IntvInstr}}
\newcommand{\oplookup}{\mathsf{lookup}}
\newcommand{\oprestrict}{\mathsf{restrict}}
\newcommand{\FinSet}{\catname{FinSet}}
\newcommand{\Bij}{\catname{Bij}}
\newcommand{\Type}{\mathbf{Type}}

\newcommand{\vpure}{\mathbf{pure}}
\newcommand{\vapp}{\mathbin{\langle * \rangle}}
% \DeclareMathOperator*{\vapp}{\langle * \rangle}
\newcommand{\fmap}{\mathbin{\langle \$ \rangle}}

\newcommand{\colorone}[1]{\textcolor{purple}{#1}}
\newcommand{\colortwo}[1]{\textcolor{blue}{#1}}
\newcommand{\ParIndex}{\mathsf{ParIndex}}
\newcommand{\Env}{\mathbf{Env}}

\newcommand{\inl}{\mathbf{inl}}
\newcommand{\inr}{\mathbf{inr}}

\newcommand{\note}{\textcolor{purple}}
\newcommand{\cur}{\mathbf{cur}}
\newcommand{\uncur}{\mathbf{uncur}}

\newcommand{\intvPt}{\mathsf{intvPt}}
\newcommand{\Endo}{\mathcal{E}}
\newcommand{\sEndo}{\mathsf{Endo}}

\newcommand{\arrow}{\rightsquigarrow}
\newcommand{\arrarr}{\mathsf{arr}}
\newcommand{\arrcomp}{\ggg}
\newcommand{\arrfirst}{\mathsf{first}}


\newcommand{\DarioMonad}{\mathsf{T}_{\mathsf{CBN}}}
\newcommand{\oprun}{\mathsf{run}}
\newcommand{\dolet}{\mathbf{let}}
\newcommand{\doletexpr}[3]{\dolet \ #1 \leftarrow #2 \ \doin \ #3}
\newcommand{\doletexprParam}[4]{\dolet_{#1} \ #2 \leftarrow #3 \ \doin \ #4}
\newcommand{\doin}{\mathbf{in}}

\newcommand{\typing}[3]{#1 \vdash #2: #3}
\newcommand{\wtyping}[4]{#1 \vdash_{#2} #3: #4}
\newcommand{\polytyping}[3]{#1 \vdash_{[-]} #2: #3}
\newcommand{\TMon}[2]{T_{#1}\left(#2 \right)}

\newcommand{\catIntv}{\mathbf{IntvStruct}}

\newcommand{\bOne}{\mathds{1}}
\newcommand{\newEndo}{\mathsf{endo}}

\newcommand{\Meas}{\mathsf{Meas}}
\newcommand{\perWorld}{\mathsf{perWorld}}
\newcommand{\valueAt}{\mathsf{valueAt}}

\newcommand{\PSh}{\mathbf{PSh}}
\newcommand{\PShBij}{\PSh(\Bij)}

\newcommand{\LangName}{\mathsf{\lambda_C}}
\newcommand{\World}{\mathsf{World}}

\newcommand{\dwtyping}[4]{#1 \mid #2 \vdash^d #3: #4}
\newcommand{\pwtyping}[5]{#1 \mid #2 \mid #3 \vdash^p #4: #5}
% \newcommand{\rtyping}[4]{#1 \vdash #2: #3 \rarrow #4}

\newcommand{\Any}{\mathsf{Any}}
\newcommand{\intvConstr}{\mathsf{intvConstr}}
\newcommand{\twin}{\mathsf{twin}}

\newcommand{\Functions}{\mathsf{Functions}}
\newcommand{\Distributions}{\mathsf{Distributions}}
\newcommand{\Constants}{\mathsf{Constants}}

\newcommand{\rone}[1]{{#1}^{(1)}}
\newcommand{\rtwo}[1]{{#1}^{(2)}}
\newcommand{\rthree}[1]{{#1}^{(3)}}

\newcommand{\Lone}{L_1}
\newcommand{\Ltwo}{L_2}
\newcommand{\Lthree}{L_3}

\newcommand{\Sbs}{\catname{Sbs}}
\newcommand{\Qbs}{\catname{Qbs}}

\newcommand{\typingone}[4]{#1 \vdash_{1}^{#4} #2: #3}
\newcommand{\etypingone}[3]{\typingone{#1}{#2}{#3}{\mathbf{e}}}
\newcommand{\ttypingone}[3]{\typingone{#1}{#2}{#3}{\mathbf{t}}}
\newcommand{\qtypingone}[3]{\typingone{#1}{#2}{#3}{\mathbf{q}}}

\newcommand{\typingtwo}[5]{#1 \mid #2 \vdash_{2}^{#5} #3: #4}
\newcommand{\etypingtwo}[4]{\typingtwo{#1}{#2}{#3}{#4}{\mathbf{e}}}
\newcommand{\ttypingtwo}[4]{\typingtwo{#1}{#2}{#3}{#4}{\mathbf{t}}}
% \newcommand{\qtypingtwo}[4]{\typingtwo{#1}{#2}{#3}{#4}{q}}

\newcommand{\etypingthree}[4]{#1 \mid #2 \vdash_{3}^{\mathbf{e}} #3: #4}
\newcommand{\ttypingthree}[6]{#1; #2 \mid #3; #4 \vdash_{3}^{\mathbf{t}} #5: #6}
% \newcommand{\qtypingtwo}[4]{\typingtwo{#1}{#2}{#3}{#4}{q}}
\newcommand{\Gammaexo}{\Gamma_{\mathsf{exo}}}
\newcommand{\Gammaendo}{\Gamma_{\mathsf{endo}}}
\newcommand{\Deltaexo}{\Delta_{\mathsf{exo}}}
\newcommand{\Deltaendo}{\Delta_{\mathsf{endo}}}

\newcommand{\typingHO}[3]{#1 \vdash_{\mathsf{HO}} #2: #3}


\newcommand{\scoringquery}[4]{\letbind{#1}{#3}{\opscore(#2)};\ \return{#4}}

\newcommand{\xletbind}[3]{\mathbf{let \ exo} \ #1 = #2 \ \mathbf{in} \ #3}
\newcommand{\nletbind}[3]{\mathbf{let \ endo} \ #1 = #2 \ \mathbf{in} \ #3}


\newcommand{\opdefault}{\mathbf{default}}
\newcommand{\opcondition}{\mathsf{condition}}
\newcommand{\opapplyIntv}{\mathsf{applyIntv}}
\newcommand{\forget}{\mathbf{forget}}
\newcommand{\semSCM}[1]{\sem{#1}_{\mathsf{SCM}}}

% Section 4
\newcommand{\IntvConstrName}{\mathbf{IntvC}}
\newcommand{\semone}[1]{\sem{#1}_{\Lone}}
\newcommand{\semtwo}[1]{\sem{#1}_{\Ltwo}}
\newcommand{\semthree}[1]{\sem{#1}_{\Lthree}}
\newcommand{\semho}[1]{\sem{#1}_{\LHO}}
\newcommand{\liftedEndo}{\widehat{\Endo}}
\newcommand{\sliftedEndo}{\sEndo}
\newcommand{\lift}{\mathsf{lift}}
\newcommand{\liftEndo}{\mathsf{endo}}
\newcommand{\catp}{\mathbf{P}}
\newcommand{\App}{\mathsf{App}}
\newcommand{\intvLabel}{\mathsf{intvLabel}}
\newcommand{\ulBij}{\underline{\Bij}}

% Section 5
\newcommand{\LHO}{L_{\mathsf{HO}}}
\newcommand{\LthreeHOcompile}[1]{\llparenthesis #1 \rrparenthesis}
\newcommand{\IntvMonadSymb}{\mathcal{I}}
\newcommand{\IntvMonad}[2]{\IntvMonadSymb_{#1}\left(#2\right)}
\newcommand{\Twinned}{\mathcal{T}}
% \newcommand{\TwinnedIntvT}[1]{\Twinned_{#1} \IntvMonadSymb}
\newcommand{\sTwinned}{\mathsf{Twinned}}
\newcommand{\sIntvMonadSymb}{\mathsf{IntvT}}
\newcommand{\sIntvMonad}[2]{\sIntvMonadSymb_{#1}\left(#2\right)}
\newcommand{\opnewEndo}{\mathsf{newEndo}}
\newcommand{\unzip}{\mathsf{unzip}}
\newcommand{\zip}{\mathsf{zip}}

% commenting
\colorlet{jyColour}{blue}

\makeatletter
\patchcmd{\@@addmarginpar}{\ifodd\c@@page}{\ifodd\c@@page\@@tempcnta\m@@ne}{}{}
\makeatother
\reversemarginpar
\newcommand{\sidenote}[3]{\marginpar{\tiny\itshape\raggedright\textcolor{#2}{\textbf{#1:} #3}}\ignorespaces} \marginparwidth=1.4cm \marginparsep=0cm
\newcommand{\jy}[1]{\sidenote{jy}{jyColour}{#1}}


% SCMs
\DeclarePairedDelimiterX{\parensbar}[1]
{(\!\delimsize\vert}
{\delimsize\vert\!)}
{#1}
\newcommand{\compileSCM}[1]{\ensuremath{\parensbar{#1}}}
\newcommand{\scmIntervene}[2]{\ensuremath \mathsf{do}_{#1, #2}}
\newcommand{\scmTwin}[1]{\ensuremath \mathsf{twin}_{#1}}
\newcommand{\IntClass}{\ensuremath{\mathsf{Intv}}}
\newcommand{\DefaultClass}{\ensuremath{\mathsf{Default}}}

\newcommand{\commentcolor}{\textcolor{codegreen}}
\newcommand{\highlightcolor}{\textcolor{blue}}
\newcommand{\SCM}{\mathcal{M}}

\title{Qualifying Dissertation for the Transfer of Status}
\author{Theo Wang}
\date{\today}

\addbibresource{refs.bib}

\begin{document}

\maketitle

% \begin{abstract}
%     Placeholder
% \end{abstract}

\section{Introduction}

% [I am interested in causality, approached it in two angles in the past year. Admittedly narrow, and misses quite a few things, that I plan to learn this year.]

% % Literature review (input by main.tex after the Introduction).
% Sections: foundations; causality as intervention; causality in
% neighbouring fields (HOQC, concurrency); synthesis.

% ---------------------------------------------------------------------
\section{Semantic foundations}\label{sec:foundations}

\subsection{Monads, algebraic theories, and their combinations}

Moggi's observation that notions of computation are monads \cite{MOGGI199155,moggi1991monadicmetalanguage} and Plotkin and Power's refinement that most of them are presented by algebraic theories of operations and equations \cite{plotkin_power_2002,PlotkinPower04CompEffectsOverview} are the foundation of everything below. An algebraic theory determines a monad, and conversely a finitary monad is a Lawvere theory \cite{lawveretheories_monads_HYLAND2007437}; the equational presentation is what makes program equivalence a matter of derivability. Hyland, Plotkin and Power \cite{sums_tensors_HYLAND200670} showed that the standard ways of combining theories have computational meaning: the \emph{sum} adjoins the operations of two theories with no interaction, so that the order of effects from the two theories is observable, while the \emph{tensor} additionally imposes that operations of one theory commute with those of the other, so that their relative order is unobservable. In the premonoidal semantics of effects \cite{PowerRobinson97premonoidal,fuhrmann1999premonoidal}, the morphisms that commute with everything are the \emph{central} ones; centrality is the established algebraic form of ``cannot be observed to happen before or after'', and so is already a primitive statement of causal independence. Handlers \cite{bauer2019algebraicalgebraiceffectshandlers} give the programming counterpart of models; comodels and runners \cite{model_comodel_PLOTKIN2008295,runners_10.1007/978-3-030-44914-8_2} describe the environments in which effectful programs execute. Operations that delimit a scope, such as catching an exception or running a nondeterministic sub-computation once, are not algebraic in Plotkin and Power's sense; Piróg, Schrijvers, Wu and Jaskelioff's \emph{scoped effects} \cite{pirog2018scoped} extend the framework to them, which matters here because normalisation, renormalising a sub-computation after conditioning, is an operation of exactly this kind (\S\ref{sec:prop-graph}).

Ordinary algebraic theories cannot express operations that create and bind \emph{names} or \emph{resources}. Staton's \emph{parameterised algebraic theories} (PATs) \cite{statonlics13,statonpopl15} extend the framework with parameters: operations consume parameters, classically choose a branch, and bind fresh parameters in each branch. Stark's free model of $\pi$-calculus name creation \cite{stark:freamp} is an earlier instance of the same discipline, and Staton used PATs to give complete equational theories of local state \cite{statonlics13} and of qubit quantum computation \cite{statonpopl15}. Kammar, Liell-Cock, Lindley, Matache and Staton \cite{kammar2026threads} have given a complete equational axiomatisation of dynamic threads (fork, wait, stop, atomic actions) as a parameterised algebraic theory, whose closed terms correspond exactly to pomsets; I return to this in \S\ref{sec:concurrency}. PATs are the formalism in which the proposal's algebraic theory of sampling, intervention and normalisation is to be stated (\S\ref{sec:prop-graph}), with intervention labels as parameters.

\subsection{Semantics of probabilistic programs}

Probabilistic programming gives a compositional syntax for statistical models and performs Bayesian inference on them compositionally \cite{vandemeent2021introductionprobabilisticprogramming}. Its denotational semantics is monadic: the Giry monad on measurable spaces \cite{kock72,semantics_ppl_staton2016}, and, for higher-order programs, the probability monad on quasi-Borel spaces \cite{Heunen_2017}. Recursion, which the proposal needs for one reading of cyclic models (\S\ref{sec:prop-expressivity}), is handled by Vákár, Kammar and Staton's domain theory for statistical probabilistic programming \cite{vakar2019domain}, in which quasi-Borel spaces are combined with $\omega$-cpos so that fixpoints and sampling coexist; probabilistic coherence spaces \cite{ehrhard2014pcoh} are the other standard model of probabilistic recursion, and the reference point for full abstraction. Implementations such as LazyPPL \cite{lazyppl}, MonadBayes \cite{monadbayes1,monadbayes2} and GenJAX \cite{becker26genjax} realise the monadic structure directly, LazyPPL through Haskell's laziness, which lets infinite-dimensional models be sampled on demand. For inference at scale, the Python systems Pyro \cite{bingham2018pyro}, on PyTorch, and NumPyro \cite{phan2019composable} and GenJAX \cite{becker26genjax}, on JAX, \emph{vectorise}: independent sub-computations and particles are batched into array operations, so that an inference algorithm is a single traced program; this is the engine the proposal intends to sit on (\S\ref{sec:prop-practice}). Conditioning itself has an algebraic account: Stein and Staton \cite{stein2021exact} treat exact conditioning as an effect with a compositional semantics in Markov categories with conditionals.

The synthetic approach of \emph{Markov categories} \cite{fritz2020synthetic,fritz2019probability} axiomatises the categorical structure needed to talk about probability without measure theory: a symmetric monoidal category with a commutative comonoid on every object, where discarding is natural (the unit is terminal) but copying need not be. In Markov categories with conditionals, disintegration and Bayesian inversion are expressed synthetically \cite{cho2019disintegration,fritz2020synthetic}, and Fritz and Klingler \cite{fritz2023dseparation} prove the soundness of d-separation, and the equivalence of the local and global Markov properties, synthetically; completeness of d-separation in the sense of Geiger, Verma and Pearl is the classical result. The \emph{representable} Markov categories of Fritz, Gonda, Perrone and Rischel \cite{fritz2023representable} are those that arise as Kleisli categories of affine commutative monads, which is the bridge between the synthetic presentation and the monadic one that our languages are interpreted in. This is the \emph{factorisation} face of causality: a Bayesian network is a string diagram in a Markov category, and the conditional independences it entails are read off the diagram.

The second piece of structure the semantics relies on is \emph{grading}. Graded monads \cite{katsumata2014_graded_strong_monad,orchard2020unifying,fujii2016towards,dorsch2019graded} and, more generally, locally graded categories \cite{levy2019locallygraded,sarkis2025string} track an index---an effect annotation, a trace, a resource---through composition. Milius, Pattinson and Schröder \cite{milius2015gradedtrace} use graded monads for trace semantics; Liell-Cock and Staton's $\mathbf{ImP}$ construction \cite{liell2025compositional} uses grading to make imprecise probability compositional. The interventional semantics of \S\ref{sec:intv-wherewestand} is built on an analogous construction in which the grade records which interventions are available.

\emph{What this gives us and what it lacks.} Markov categories give factorisation but no primitive notion of intervention; the diagrammatic causal-model literature (\S\ref{sec:diagrammatic}) adds it by surgery on diagrams. Neither offers a notion of program \emph{context}, which is what contextual equivalence needs; and the monadic semantics in which our languages are interpreted is stated for particular categories rather than synthetically (\S\ref{sec:prop-synthetic}).

% ---------------------------------------------------------------------
\section{Causality as intervention}\label{sec:intv}

\subsection{Structural causal models and Pearl's hierarchy}

Pearl's structural causal models (SCMs) \cite{pearl2009causality,Pearl2018_book_of_why} are systems of equations $X_i = f_i(\mathrm{Pa}_i, U_i)$ in which each endogenous variable is a deterministic function of its parents and of an exogenous noise variable; a prior on the noise makes an SCM a probabilistic model, and the equations make it more than a distribution. The surplus is organised by the \emph{causal hierarchy} \cite{bareinboim:etal20}: rung 1 asks observational questions; rung 2 asks \emph{interventional} questions, answered by replacing the equation for $X$ with a constant (the $\dok$ operator); rung 3 asks \emph{counterfactual} questions about a particular run, answered by abduction, action and prediction, implemented graphically by \emph{twinning}, which duplicates the model into worlds that share the noise \cite{Balke1994}. Both operations are \emph{syntactic transformations of the model}. The Neyman--Rubin potential-outcomes framework \cite{rubin1974} is the other dominant formalism; its correspondence with SCMs is well understood \cite{ibeling2023comparing}. Richardson and Robins' single-world intervention graphs \cite{richardson2013swig} present an intervention by splitting a node into its natural and its fixed value, which is the graphical counterpart of the intervention labels of our languages; they are deliberately single-world, and cross-world quantities such as natural direct effects need further assumptions there. Interventions need not be perfect: the literature on stochastic and soft interventions \cite{dawid2002_stoch_intv,didelez2006_stoch_mech_intv,Eberhardt_Scheines_2007,levis2024stochasticinterventionssensitivityanalysis} replaces an equation by another mechanism rather than a constant.

\paragraph{Graphs, Markov conditions and fine tuning.} The graph of an SCM is derived from its equations: an arrow $X\to Y$ when $f_Y$ reads $X$. The \emph{causal Markov condition} says that the induced distribution factorises along the graph, and \emph{d-separation} \cite{pearl2009causality,spirtes2000causation} reads off the conditional independences that this forces; the \emph{faithfulness} assumption of causal discovery says that there are no others. Faithfulness can fail, and Barrett, Lorenz and Oreshkov \cite{barrett2019quantumcausalmodels} give (their \S2.4) the two-line example that the proposal returns to: $Y = X + \lambda_Y \bmod 2$ with $\lambda_Y$ uniform. Here $X$ is a direct cause of $Y$, since $f_Y$ reads $X$, yet $P(Y\mid \dok(X)) = P(Y)$: no intervention on $X$ is detectable at $Y$, and no signal can be sent from $X$ to $Y$. The independence is \emph{fine tuned}: it holds for one distribution on $\lambda_Y$ and for no other. The distinction between independences forced by the graph and accidental ones is thus the distinction between an intensional datum (the graph) and an extensional one (what interventions reveal), and it is the same distinction as the one between the dataflow of a program and its contextual equivalence class (\S\ref{sec:intv-wherewestand}).

\paragraph{Identifiability and the do-calculus.} When only the graph and observational data are given, the question whether an interventional quantity $P(y\mid\dok(x))$ is determined is \emph{identifiability}. Pearl's \emph{do-calculus} \cite{pearl1995docalculus} is a set of three inference rules, each licensed by a d-separation condition in a graph with some arrows cut, that rewrite interventional expressions into observational ones; Huang and Valtorta \cite{huang2006complete} and Shpitser and Pearl \cite{shpitser2006identification,JMLR:v9:shpitser08a} prove it complete for identification, with an algorithm that either produces a formula or a counterexample. The do-calculus is therefore an \emph{equational theory} of interventional queries indexed by a graph, and the natural object to compare with the program-level equational theories of \S\ref{sec:intv-wherewestand} (\S\ref{sec:prop-graph}). Causal discovery \cite{spirtes2000causation,peters_janzing_schoelkopf_2017_book}, which infers the graph itself from data, remains outside the scope of this thesis.

\paragraph{Cyclic models.} Pearl's models are acyclic. Bongers, Forré, Peters and Mooij \cite{Bongers_2021} extend SCMs to \emph{cyclic} systems of equations and give them meaning through \emph{solution functions}: for a fixed value of the noise, a solution is an assignment satisfying all equations at once, and the model defines a distribution when a solution exists and is unique (or measurably selectable). Forré and Mooij \cite{forre2017markov} replace d-separation by $\sigma$-separation and prove the corresponding Markov properties for cyclic graphs, where the acyclic ones fail. Mooij, Janzing and Schölkopf \cite{mooij2013ode} show where such models come from: a system of ordinary differential equations that converges to an equilibrium induces a cyclic SCM whose solutions are the equilibria, and interventions correspond to clamping a variable. This is one reading of feedback, in which the randomness is drawn once and the cycle is a fixpoint of deterministic equations. The other reading, standard in time-series and dynamic Bayesian networks, unrolls the cycle in time: each iteration draws fresh randomness and the model is an acyclic model over indexed copies of the variables, a stochastic process. The two agree in the deterministic equilibrium case and differ in general. Neither has a programming-language account, in which feedback would be a loop or a recursive definition (\S\ref{sec:prop-expressivity}).

\emph{What this gives us.} The causal structure of an SCM is its dataflow: which variables each equation reads. Interventions probe it, but, by fine tuning, do not determine it. The semantic question is when two models have the same causal structure, and the answer must say which of the two one means.

\subsection{Equivalences of causal models and logical accounts}

Bongers, Forré, Peters and Mooij \cite{Bongers_2021} give the principal account of equivalence for SCMs at all three rungs: two models are observationally, interventionally or counterfactually equivalent when they agree on the corresponding class of queries, with the counterfactual case defined through the two-world twinned model, and the treatment covers cyclic models through their solution functions. At rung 2, several further notions exist. Pearl's causal Bayesian networks are families of interventional distributions compatible with a graph. \emph{Causal spaces} \cite{park2024measuretheoreticaxiomatisationcausality,buchholz2024products} axiomatise rung 2 measure-theoretically: a causal space is a probability space together with a family of \emph{causal kernels}, one for each sub-$\sigma$-algebra, giving the distribution after intervening on the events it measures, subject to consistency axioms; products of causal spaces are defined, and the construction makes no reference to a graph or to mechanisms. \emph{Causal abstraction} \cite{rubenstein2017consistency,beckers2019abstracting,rischel2021abstraction,lorenztull2026abstraction,englberger2025causalabstractionscategoricallyunified} studies maps between models of different granularity, a coarse model being an abstraction of a fine one when interventions on the coarse variables and their images on the fine ones commute with the map; Rischel and Weichwald \cite{rischel2021abstraction} and the later string-diagram line make this a morphism of models in a category. At the level of graphs rather than models, Markov equivalence \cite{verma1990equivalence} is what discovery algorithms can recover. On the logical side, Ibeling and Icard \cite{ibeling2020hierarchy} characterise the three rungs by three probabilistic logics of strictly increasing expressivity whose semantics is defined on SCMs through Pearl's operations, and Icard has proposed probabilistic Turing machines as the computational counterpart of SCMs \cite{icard2017programs,ibeling2019openuniverse}.

\emph{What this gives us and what it lacks.} These are all extensional characterisations of interventional structure, and they are the yardstick against which our contextual equivalences are measured (\S\ref{sec:intv-wherewestand}). None of them is formulated as a contextual equivalence of programs, and none treats higher-order models; the proposal asks how each is related to the graded semantics (\S\ref{sec:prop-relations}).

\subsection{Diagrammatic and categorical causal models}\label{sec:diagrammatic}

Fong's \emph{causal theories} \cite{fong2012causaltheories} present a Bayesian network as a symmetric monoidal category freely generated by its graph, with models as functors into stochastic maps; this is the first place where the graph is the \emph{primitive} and the model is derived from it. Jacobs, Kissinger and Zanasi \cite{jacobs2019causal} present interventions as \emph{string-diagram surgery} in this setting, and introduce \emph{comb disintegration}: the use of combs, processes with a hole, to express the classical causal-inference computations, which is the first appearance of the comb notion of \S\ref{sec:hoqc} in classical causal inference and a diagrammatic form of identification. Lorenz and Tull \cite{lorenz2023causalmodelsstringdiagrams} develop causal models in string diagrams systematically, including open models with inputs, general interventions, counterfactuals, and the identifiability of causal effects treated diagrammatically; a later line presents causal abstraction as natural transformations between such models \cite{lorenztull2026abstraction,englberger2025causalabstractionscategoricallyunified}. Bayesian inference in this style is handled by \emph{normalisation boxes}, which renormalise locally inside a diagram \cite{jacobsszelesstein2025inference}, and Agrawal et al.\ \cite{agrawal2026equational} obtain equational principles for causal Bayesian networks from them. Zhang et al.\ \cite{zhang2026semanticseparation} characterise d-separation as a noninterference-style separation property, which makes explicit a connection implicit in all of this work and in \S\ref{sec:intro}: invariance under interventions is a noninterference property \cite{goguen1982noninterference}, and type systems for noninterference, from the dependency core calculus onward \cite{abadi1999dcc}, are the programming-language way of carrying a dependency graph in the \emph{types} rather than in the semantics.

\emph{What this gives us and what it lacks.} This is the closest existing bridge between the factorisation face and the intervention face of causality, and it is compositional in the functorial sense. Its equational theories are, however, stated per graph, and it has no notion of program context or of higher-order models. It is where the graph-centred direction of the proposal (\S\ref{sec:prop-graph}) starts: the surgery, comb and normalisation-box calculi are the diagrammatic form of the do-calculus, and an algebraic theory in which sampling, intervention and normalisation are all operations is the programming-language form that does not yet exist.

\subsection{Causal probabilistic programming languages}

Causal probabilistic programming aims to make intervention and twinning primitives of an expressive probabilistic language. Libraries such as DoWhy \cite{pywhy,blobaum2024dowhygcm} and the $\dok$ operators of Pyro and PyMC \cite{bingham2018pyro,abrilpla2023pymc} manipulate models as graphs or computation graphs and stop at rung 2 (DoWhy's GCM module reaches rung 3 on fitted mechanisms). Witty et al.\ \cite{witty2019synthesis,witty2023bayesian} treat an intervention as a source transformation that replaces subterms, which is powerful but non-compositional; Brulé's Whittemore \cite{brule2018whittemore} is an early embedded causal DSL.

Four systems reach rung 3 natively. OmegaC \cite{tavares2021language} obtains interventions by a lazy evaluation strategy that retroactively modifies closures, reminiscent of dynamic binding, and counterfactuals by making every program a deterministic function of explicit noise. ChiRho \cite{chirho} implements both as effect handlers over Pyro, with worlds realised by tensor broadcasting along a batch dimension and each intervention doubling the world count, so that its relation to Pearl's twinning is not immediate. MultiVerse \cite{perov2020multiverse} uses trace inspection and re-execution. Kiesel et al.\ \cite{kiesel2023whatif} answer counterfactual queries on ProbLog programs by a transformation proved correct against Pearl's semantics.

\emph{What this gives us and what it lacks.} These systems show that the primitives are implementable at higher order, by non-standard evaluation techniques, and ChiRho shows that they can be vectorised. What they lack is a denotational and equational foundation: a semantics in which intervention is a mathematical operation rather than an evaluation strategy, and an account of when two causal programs are equivalent. That is what my first-year work supplies; what it does not yet supply is a semantics \emph{of} these systems, or an implementation that runs at their scale (\S\ref{sec:prop-practice}, \S\ref{sec:prop-relations}).

\subsection{Where I stand}\label{sec:intv-wherewestand}

Our paper, to appear at POPL 2027 \cite{wang2026causalppl}, with Dario Stein (joint first author), Michael Lee, Eli Bingham, John Feser, Ohad Kammar, Sam Staton and Jeremy Yallop (an early version appeared at the LAFI workshop \cite{lafi-wang}; artifact \cite{ARTIFACT_wang_2026_23207221}), reconstructs Pearl's hierarchy as a hierarchy of three first-order probabilistic languages. $\Lk{1}$ is a standard first-order probabilistic language with conditioning by a final score. $\Lk{2}$ exposes hooks in the dataflow as typed free variables, the \emph{intervention labels}, read through $\mathsf{applyIntv}$; an intervention $\dok(\ell\mapsto v)\{t\}$ is substitution of $v$ for $\ell$. $\Lk{3}$ annotates each binding as exogenous or endogenous, and \emph{twinning} $\mathsf{twin}_n$ is a vectorisation transformation from $\Lk{3}$ to $\Lk{2}$ that duplicates the endogenous bindings into $n$ worlds sharing the exogenous ones. Program equivalence at rung $k$ is \emph{contextual} with respect to rung-1 denotational equality, in Milner's sense: two programs are $\equiv_k$ when they are denotationally equal in every rung-$k$ context, and varying the admissible interventions and world counts gives a spectrum of equivalences. For representable acyclic SCMs compiled into $\Lk{3}$, the observational, interventional and counterfactual equivalences of Bongers et al.\ coincide with the equivalences $\equiv_3^{\mathrm{obs}}$, $\equiv_3^{\mathrm{intv}}$ and $\equiv_3^{\text{2-world}}$ of the compiled programs when the model is concrete, the forward implication holding in general; the full $\equiv_3$ quantifies over all $n$-world twinnings, and whether it is strictly finer than $\equiv_3^{\text{2-world}}$ is open. The denotational semantics interprets interventions by \emph{grading}, in a construction reminiscent of $\mathbf{ImP}$ \cite{liell2025compositional}, and twinning by world-polymorphic vectorisation. It is adequate for the contextual equivalences unconditionally and fully abstract when the underlying model is \emph{concrete}, meaning that values correspond to global elements and global elements separate morphisms. The construction is parametric in an ``$\Lk{1}$-model'': a distributive category with a commutative, affine probability monad and enough structure to interpret the base language, with cartesian closure added for the higher-order language; $\mathbf{Set}$, standard Borel spaces and quasi-Borel spaces are the instances, and all three are concrete. The denotation then motivates a higher-order language $\LHO$ with cross-world interventions, into which $\Lk{3}$ compiles in a semantics-preserving way, and which has a definitional Haskell implementation. \fillin{one sentence on which parts are yours: e.g.\ the contextual-equivalence framework and the full-abstraction proofs}

Two consequences frame the proposal. First, $\equiv_2$ is the extensional notion: it identifies exactly the programs that no interventional observer can tell apart. The fine-tuned model of Barrett, Lorenz and Oreshkov is a one-line program, $\mathtt{let}\ \lambda = \mathtt{bernoulli}\ 0.5\ \mathtt{in}\ \mathtt{return}\ (\mathtt{xor}\ t'\ \lambda)$ with $t' = \mathtt{applyIntv}\ \ell\ t$, and it is $\equiv_2$ to the program that never reads $t'$; the dataflow, which is the graph, is finer than $\equiv_2$, and the full-abstraction theorem is a theorem about the extensional notion only. Second, the limitations relevant to the proposal are these: the languages are acyclic; the relation between the contextual equational theory and the graph-indexed equational theories of \S\ref{sec:diagrammatic}, the do-calculus included, is not worked out; the categorical assumptions are a list of properties of particular categories rather than a synthetic axiomatisation; the first-order semantics is over normalised probability, while $\LHO$ widens to the unnormalised measure monad to internalise conditioning, so that programs with observations denote unnormalised, weighted maps; and the models are measure-theoretic rather than finite-dimensional, which matters for the signalling face (\S\ref{sec:synthesis}).

% ---------------------------------------------------------------------
\section{Causality in neighbouring fields}\label{sec:order}

The two fields of this section are positioned, not surveyed: my results are not in them, I know them at the level of their main definitions and results, and their detailed study is a milestone of the plan (\S\ref{sec:plan}). They are included because each has a notion of causality that the proposal's fourth direction (\S\ref{sec:prop-other}) wants to relate to the interventional one. The one item of that direction that is a theorem rather than a position, the classical bridge D2, needs only Kissinger and Uijlen's first-order and one-way-signalling types and their classical instance over $\Mat$, which I do know.

\subsection{Higher-order quantum processes and signalling}\label{sec:hoqc}

In the standard semantics of quantum computation \cite{QPL_SELINGER_2004,PQP_Coecke_Kissinger_2017}, types denote finite-dimensional spaces of density matrices and programs denote completely positive maps; the deterministic processes, those that preserve the discarding map, are the trace-preserving ones, and the category $\CP$ of all completely positive maps is compact closed, which the subcategory of deterministic ones is not. Higher-order quantum computation studies transformations of such processes rather than of states. Chiribella, D'Ariano and Perinotti's \emph{quantum combs} \cite{quantum_circuit_architecture_Chiribella_2008,chiribella2009networks} are processes with holes in a fixed causal order; their probabilistic counterparts, \emph{testers}, and Gutoski and Watrous's \emph{strategies} \cite{gutoski2007games} describe multi-round interactions with sub-normalised outcomes, and the general theory of supermaps forms a hierarchy of types \cite{bisio2019hoqt}. Oreshkov, Costa and Brukner's \emph{process matrices} \cite{oreshkov2012noncausal} are the most general bipartite process operators compatible with quantum mechanics holding locally in each party's laboratory; the multipartite theory and the notion of \emph{causal separability}, in which the order of later parties may depend classically on earlier outcomes, are due to Araújo et al.\ and Oreshkov and Giarmatzi \cite{araujo2015witnessing,oreshkov2016separable}. Some process matrices are causally non-separable; the \emph{quantum switch} \cite{chiribella2013switch}, which applies two operations in an order controlled by a qubit in superposition, is the standard example \cite{araujo2015witnessing}, and which indefinite-order processes can be built as circuits with classical or quantum control of order is the subject of Wechs et al.\ \cite{wechs2021qcqc}.

Categorically, Kissinger and Uijlen's $\CausC$ construction \cite{kissinger2019categorical} builds, over a \emph{precausal} category---compact closed, with a chosen discarding structure and further axioms, the examples being $\CP$ on finite-dimensional systems and $\Mat$, finite matrices of non-negative reals---a $*$-autonomous category whose types are signalling constraints: first-order processes, non-signalling processes, one-way signalling processes, and so on. The classical example matters for this thesis: every construction of the quantum theory goes through for $\Mat$, where first-order processes are stochastic maps, so that combs, testers and signalling types have a purely classical meaning. The dual of the non-signalling type contains all bipartite process matrices, causally separable or not; separability is a property inside that type, not a type, and Simmons shows, as \emph{affine universality} of definite orders, that every process matrix is an affine combination of definite-order processes \cite{simmons2024thesis}. Simmons's thesis \cite{simmons2024thesis} is the fullest development of this line. It rebuilds the construction over \emph{additive} precausal categories, which gives $\CausC$ products, coproducts and probabilistic choice as type operators; it adds a non-commutative sequential operator for one-way signalling, under which $\CausC$ models Guglielmi's BV-logic \cite{simmons2022bv}; it defines $n$-ary \emph{graph types} encoding compatibility with an arbitrary DAG of signalling relations, which is the one existing place where a causal graph is a \emph{type}; it surveys candidate ways of admitting sub-normalised, partial processes into the construction and where each falls short; and it introduces \emph{causal logic}, pomset logic with directed axiom links, whose proof nets characterise exactly the causally consistent wirings of black boxes \cite{simmons2024thesis,simmons2024completelogic}. Simmons's finite-causal-basis axiom excludes infinite-dimensional and measure-theoretic bases, which is one of the two obstructions discussed in \S\ref{sec:synthesis}.

\emph{Quantum causal models} transfer Pearl's framework to quantum processes: Costa and Shrapnel \cite{costa2016quantumcausalmodelling} gave the first process-matrix-based definition, Allen et al.\ \cite{allen2017quantumcommoncauses} the common-cause line, and Barrett, Lorenz and Oreshkov \cite{barrett2019quantumcausalmodels} the general framework, in which a node is an input--output pair, the process operator factorises as a product of commuting operators (a quantum Markov condition), interventions are local instruments at nodes, d-separation in the graph is sound and complete for a quantum analogue of conditional independence, and the three rules of the do-calculus have quantum generalisations. Three features of their framework (their \S\S2.4, 4.1 and 4.5) are relevant here. Causal structure is defined for a \emph{unitary} realisation by a no-influence condition between inputs and outputs, and the DAG of a model is derived from it; the DAG is in general a strict subgraph of the \emph{connectivity} DAG of the circuit, since a unitary may have no-influence relations that its wiring does not show; and, as in the classical case, a direct cause need not permit signalling in the marginal process, the signalling being washed out by a fine-tuned state on an unobserved system. There are thus three strictly nested notions, wiring, causal structure and signalling, of which the middle one is the primitive. The cyclic extension \cite{barrett2021cyclic} gives a cyclic-causal-model account of some processes without definite order, including the switch, which is the quantum form of the observation that feedback and indefiniteness of order are related.

\emph{Position.} Two things follow for this thesis. First, signalling types are the extensional face of causal structure in the same way that interventional equivalence is the extensional face of dataflow, and the fine-tuning phenomenon is the same in both: a signalling constraint says that a local intervention at one node is undetectable at another, which is the interventional notion stated for processes with slots. Second, quantum causal models are the one existing place where factorisation, intervention and signalling are all present, and they are second-order processes with nodes as slots, exactly as a program with intervention labels is a process with slots; that is why they are a reading milestone and the starting point of \S\ref{sec:prop-hoqc}.

\subsection{Order in concurrency theory}\label{sec:concurrency}

Concurrency theory has long modelled causal order explicitly. Lamport's happens-before relation \cite{lamport1978time} is the partial order an observer of a distributed system can reconstruct. Mazurkiewicz traces \cite{diekert1995traces} are words modulo the commutation of independent letters, the cheapest partial-order model. Pratt's pomsets \cite{pratt1986pomsets}---partially ordered multisets of events, up to isomorphism, with series and parallel composition---generalise words to partially ordered words, and Gischer \cite{gischer1988pomsets} axiomatised the equational theory of series-parallel pomsets. Winskel's event structures \cite{winskel1986eventstructures,nielsen1981petri,winskel1995models} add a conflict relation to the causal order, so that a single structure records both what must precede what and what excludes what, and give a domain-theoretic account of nondeterministic concurrent computation. Concurrent Kleene algebra \cite{hoare2011cka} adds a commutative parallel operator to Kleene algebra, related to sequential composition by an exchange law; Kappé, Brunet, Silva and Zanasi \cite{kappe2018cka} prove that series-parallel pomset languages closed under Gischer's subsumption order are its free model, building on Laurence and Struth's completeness theorem for bi-Kleene algebra \cite{laurence2014completeness}. Within the algebraic-effects tradition, Dvir, Kammar and Lahav \cite{dvir2022sharedstate} give a monadic semantics of shared-state concurrency from an equational theory, and Kammar et al.\ \cite{kammar2026threads} axiomatise dynamic threads as a parameterised algebraic theory whose closed terms are exactly pomsets, with a completeness theorem. Clairambault, de Visme and Winskel \cite{clairambault2019quantumgames} give an event-structure game semantics for quantum programs.

\emph{Position.} In these models the causal graph is \emph{explicit}: a pomset or event structure is a graph of events, and the semantics is the graph. This is the opposite extreme from the interventional line, where the graph is implicit in the dataflow and only its extensional shadow is characterised. Whether the interventional notion has a counterpart in these models, intervening on an event, and whether the signalling types of \S\ref{sec:hoqc} are the extensional reading of a pomset, are the open-ended questions of \S\ref{sec:prop-concurrency}.

% ---------------------------------------------------------------------
\section{Synthesis: three notions of causality and the gap}\label{sec:synthesis}

Table \ref{tab:synthesis} summarises the review. Each column is one extensional notion of causality; each row asks the same question of it.

\begin{table}[ht]
\centering\small
\begin{tabularx}{\textwidth}{>{\raggedright\arraybackslash}p{2.6cm}>{\raggedright\arraybackslash}X>{\raggedright\arraybackslash}X>{\raggedright\arraybackslash}X}
\toprule
 & \textbf{Intervention} & \textbf{Factorisation} & \textbf{Signalling}\\
\midrule
Intensional datum & dataflow of a program / SCM & DAG & wiring of a circuit; comb order; pomset\\
Extensional characterisation & contextual equivalence under interventions & conditional independence / d-separation & signalling types ($\CausC$); causal (non-)separability\\
Gap between the two & fine tuning: dataflow finer than $\equiv_2$ & faithfulness: graph finer than the independences & fine tuning: causal structure finer than signalling \cite{barrett2019quantumcausalmodels}\\
Correspondence theorem & full abstraction for concrete models (\S\ref{sec:intv-wherewestand}) & soundness of d-separation, synthetically \cite{fritz2023dseparation}; completeness classically & types of $\CausC$ characterise signalling; no programming-language counterpart\\
Equational theory & contextual equivalences $\equiv_k$ & do-calculus; normalisation-box equations & BV and causal logic\\
Semantics of programs & first-order and higher-order, acyclic & Markov categories; interventions by surgery & none for signalling; pomsets for concurrent programs\\
Next hard case & cyclic SCMs & cyclic and quantum causal models & causally non-separable processes; conflict\\
\bottomrule
\end{tabularx}
\caption{Three notions of causality and where the first-year work stands.}\label{tab:synthesis}
\end{table}

Four remarks. First, the asymmetry: for intervention the correspondence between intensional and extensional structure is a theorem, and the intensional datum is a program; for signalling the extensional characterisation exists, as a type system, but there is no programming language whose syntax it characterises; for concurrency the intensional datum is explicit and there is no interventional reading of it. Second, the row ``gap between the two'' is the same phenomenon in all three columns, and it is what motivates the graph-centred direction of the proposal: in each column the intensional datum is strictly finer than what any observer can recover, and the literature's response in each is to take the graph as primitive, as a Markov condition, a no-influence condition, or a type. Third, non-DAG structure is the next hard case in every column, but it is not the same case: feedback in cyclic SCMs, causal non-separability in quantum processes, and conflict in event structures are three different departures from a fixed order, and a unifying account must distinguish them. Fourth, and this is the gap: there is no framework in which the three columns are structure on one category. Markov categories cover factorisation and, through surgery, part of intervention; $\CausC$ covers signalling for quantum and finite classical processes; quantum causal models relate all three, with a compositional presentation in Lorenz and Tull's string diagrams, but are not formulated as a semantics of programs. The graded semantics of \S\ref{sec:intv-wherewestand} covers the first column; it is a family of morphisms indexed by an intensional datum, the grade, with a reindexing action, which is the common structure a generalisation would start from.

Two obstructions should be displayed rather than hidden. The first is \emph{finite dimension}: the $\CausC$ construction needs a compact closed base with invertible dimensions, which $\Mat$ and finite-dimensional $\CP$ provide, but measurable spaces and quasi-Borel spaces do not. Any bridge from the interventional semantics to signalling types is therefore stated on the finite fragment. The second is \emph{sub-normalisation}: conditioning in $\LHO$ denotes unnormalised, weighted maps, which do not land in a first-order $\CausC$ type as it stands; testers and strategies \cite{chiribella2009networks,gutoski2007games} accommodate sub-normalised families only, so a renormalisation is needed even to reach that setting, Simmons \cite{simmons2024thesis} surveys the ways of admitting partial processes into $\CausC$ and finds each wanting in some respect, and normalisation boxes \cite{jacobsszelesstein2025inference} are the diagrammatic tool for renormalising after conditioning.


\section{Literature Review}

\subsection{Foundations of Probabilistic Programming}

\paragraph{From Bayesian Networks to Probabilistic Programming} A central tool in Bayesian probabilistic modelling is Bayesian networks \todo{cite}, a graphical representation of a joint distribution. Formally, a Bayesian network over a certain set of variables $V$ is a directed acyclic graph over $V$, a domain $D_x$ for each variable $x \in V$, and for each variable $v \in V$, a conditional distribution $P(x \mid \mathsf{Parents}(x))$. This data can be seen as a Markov factorisation of a joint distribution $$P(x_1...x_n) = \prod_{x\in V} P(x \mid \mathsf{Parents}(x))$$
Bayesian networks are useful because their graphical structure encodes the conditional independencies between variables (a variable is assumed to only depend on its parents), exploited to make Bayesian inference (computing conditional distributions) more efficient. \todo{cite again} 

For instance consider the scenario where we are trying to model a car, which has fuel with probability $0.8$, battery with probability $0.9$, and only starts if it has both fuel and battery. Then we can model the joint distribution between $\mathsf{fuel}$, $\mathsf{battery}$ and $\mathsf{carStarts}$ with the following Bayesian network: 

\todo{Draw Bayesian Network}

Probabilistic Programming languages can be seen as a compositional generalisation of Bayesian networks: they provide a compositional syntax for specifying probabilistic models and performing inference on them. Like for Bayesian networks, we can write program represent joint distributions, like in \todo{}. Unlike Bayesian networks, we can also record observations like $\mathsf{carStarts}=\true$ using the $\opcondition: \booleans \to 1$ operation, like in \todo{}, or more general `soft' conditioning using $\opscore: \reals^+ \to 1$. \todo{Describe 2 conditionings Operationally}. Finally, general inference algorithms like \todo{} allows one to sample from the conditional distribution that the program represents. 

\todo{Write program}

\paragraph{Semantics of PPLs} The semantics of probabilistic programming languages are generally given in a category $\catc$ that `supports probability theory' with a normalised probability monad $\Prob$. In the discrete case, one considers $\Set$ and the finitely supported distributions monad. Because $\Set$ is cartesian-closed, one can not only interpret first order programs, but also higher order ones like $(X \to \Prob(Y)) \to \Prob(Z)$. In the more general continuous case, one often uses categories like the category of measurable spaces $\Meas$ (or subcategories thereof) with the Giry monad $G: \Meas \to \Meas$ \todo{cite}. Because $\Meas$ is famously not cartesian-closed, several works \cite{semantics_ppl_staton2016,Heunen_2017} have developed the category of quasi-Borel spaces (QBS) as an alternative in which higher order probabilistic programming is possible. Many extensions to QBS have since been developed, such as $\omega$QBS \todo{cite}, which enables recursion, and quasi-Borel families (QBFs), which enable dependent types \todo{cite}. 

Normalised probability monads only allow one to interpret sampling; interpreting conditioning (or scoring in general) will require \emph{unnormalised} probability monads $\Dist$. An unnormalised probability monad should be a super-monad of the normalised counterpart. Without going into too much detail, there are many possible choices for $\Dist$ with different trade-offs. For instance, we can consider $\Dist(X) = \Prob(X \times \reals^+)$, which enables general scoring using a writer effect; or indeed the subdistribution monad $\Dist(X) = \Prob(X + 1)$, which only supports (hard) conditioning. \todo{Normalisation can be thought of as...}

\paragraph{Synthetic Foundations}

% \subsubsection{Monadic Probabilistic Programming}

% \begin{itemize}
%     \item Bayesian networks at the centre of model based reasoning. Nice notation for statistical models. Conditional independencies. Denotation as just a distribution. Efficient representation for computing conditional distributions, hence "Bayesian network". 
%     \item Probabilistic programming \cite{vandemeent2021introductionprobabilisticprogramming} is programming paradigm that allows the user for declaratively specifying probabilistic models, and compositionally performing Bayesian inference on those models. \todo{Cite a few languages.} Think: compositional generalisation of Bayesian networks. 
%     \item Core primitives: sample and score. + Example program of model + conditioning statement, that allows for the computation of CONDITIONAL probabilities.
%     \item Semantics usually given by a monad: Category + Distribution Monad, or unnormalised Distribution monad to allow conditioning. Discrete case, as above, is in Set + Dist monad. Cont case is more complex, with a history of giving semantics to it using first Meas (or subcats of Meas) + Giry monad. But then no function spaces. More recently, model found for QBS + Prob monad that supports higher order reasoning.
%     \item Many works step from the monadic semantics, like lazyppl + efficient lazy inference algorithm, or like Liell-Cock and Staton's ImP construction. We do sth similar to that. 
% \end{itemize}

% \subsubsection{Synthetic Foundations}

% \begin{itemize}
%     \item Markov Categories: synthetically axiomatises the categorical structure to talk about probability. 
%     \item Definition: SMC + copy-discard structure, where discard is natural. 
%     \item String diagrammatic reasoning, elegantly prove properties like X or Y. 
%     \item A Bayesian network is now just a string diagram. 
% \end{itemize}


\subsection{Foundations of Causal Modelling}

\subsubsection{Pearl's Three Rungs of Causality}

\subsubsection{Semantics of Causality}

\paragraph{Causal Models}

\paragraph{Causal Structures}

% \subsubsection{Pearl's causality, in 3 rungs}

% Directly move the explanation of the POPL paper to here.
% \begin{itemize}
%     \item Rung 1: Bayesian Networks.
%     \item Rung 2: Interventions
%     \item Rung 3: Counterfactuals
% \end{itemize}

% \subsubsection{Equivalences of causal models}
% Copy from paper and tweak.

% \subsubsection{Logical, diagrammatic, categorical accounts of causality}

% Copy from paper and tweak. Those references are exactly the ones I need to mention. 

\subsection{Causal Probabilistic Programming Languages}

\subsubsection{Summary of our POPL 2027 paper}


\paragraph{Contribution 1: semantic foundations of causal programming}
Our first contribution is to develop a semantic foundation of causal models by taking the analogy with probabilistic programs seriously. In \cref{section:hierarchylanguages}, we reconstruct Pearl's rungs of causality as a hierarchy of simple first order probabilistic programming languages $\Lone, \Ltwo, \Lthree$.  Each language exposes more intension than the previous one using syntactic annotations, and causal inference is performed using program transformations that consume these syntactic annotations: $\Lone$ is a standard first order PPL; $\Ltwo$ exposes typed holes named \emph{intervention labels}, canonically of type $\Maybe{A}$, allowing substituting to optionally perform an intervention; $\Lthree$ annotates each let binding with $\mathbf{exo}$ or $\mathbf{endo}$, so that twinning can be formulated as a vectorisation transform. The same idea underlies the intervention-as-program-transformation view of Witty~\cite{witty2021sbi,witty2023bayesian}, the logical account by Ibeling and Icard~\cite{ibeling2020hierarchy}, and the string-diagrammatic accounts of Jacobs et al.~\cite{jacobs2019causal} and Lorenz and Tull~\cite{lorenz2023causalmodelsstringdiagrams}. 

What is new is the observation that this simple setting yields a semantic characterisation of Pearl's rungs of causality (\cref{section:scms}). We develop a unified account of contextual equivalences that recovers the full spectrum of causal model equivalences by varying the class of admissible contexts, e.g.~Rung 3 ($\Lthree$) programs are Rung 3 equivalent ($\equiv_3$) if they behave the same under all twinnings and interventions; Rung 2 equivalent ($\equiv_3^{\mathsf{intv}}$) if under all interventions; Rung 1 equivalent ($\equiv_3^{\mathsf{obs}}$) if they induce the same distribution under no intervention. The subscript $3$ marks that we compare $\Lthree$ programs throughout; the superscript marks the rung of context tested against ($\mathsf{obs}$, $\mathsf{intv}$, or none for full rung-3 twinning). These equivalences form a strict inclusion hierarchy $\equiv_3 \subsetneq \equiv_3^{\mathsf{intv}} \subsetneq \equiv_3^{\mathsf{obs}}$, an analogue of the causal hierarchy's non-collapse \cite{bareinboim:etal20}, i.e.~programs indistinguishable at a lower rung may still differ at a higher one.

In \cref{subsection:relationtoscms} we connect the framework to the theory of SCMs, showing that $\Lthree$ extends SCMs conservatively: each SCM can be compiled into $\Lthree$, preserving the semantics of intervention and twinning (\cref{thm:scmcompilationcorrectness}). Through this compilation, we show that the hierarchy of SCM equivalences developed by Bongers et al.~\cite{Bongers_2021} is naturally subsumed by our program equivalences (\cref{thm:hierarchyBongers}).

\paragraph{Contribution 2: Typed abstractions and denotational semantics for expressive causal probabilistic programming}
Our second contribution is the development of $\LHO$, a semantically-founded, expressive, and typed language for higher-order causal probabilistic programming. $\LHO$ has a similar programming model to ChiRho \cite{chirho,lafi-wang}: it is expressive as a programming language (higher order, with control flow) and as a causal calculus (rich, higher order interventions, see \Cref{fig:intvkinds}); we illustrate it with examples such as mediation analysis, which decomposes a causal effect into a direct path and a path through an intermediate variable (\Cref{fig:mediationanalysis}). However, unlike ChiRho, it extends the first-order languages in a semantic-preserving way (\cref{thm:L3LHOsemanticpreservation}). 

Rung 2: To manage interventions, we track intervention labels using \emph{grading} \cite{orchard2020unifying,katsumata2014_graded_strong_monad,dorsch2019graded,fujii2016towards,sarkis2025string}, similarly to the $\mathbf{ImP}$ construction in \cite{liell2025compositional} or intervention points in \cite{lafi-wang}. We define a graded reader monad $\IntvMonad{[\tau_1...\tau_n]}{X} \eqdef \prod_i\tau_i \to \Meas(X)$ graded by lists of intervention types $[\tau_1...\tau_n]$, where $\Meas$ is the unnormalised probability monad. A causal model of type $\IntvMonad{[\tau_1...\tau_n]}{X}$ is an intervenable probabilistic program with $n$ typed intervenable holes, each of type $\tau_i$, and a $\opdo$ intervention on the $i$th hole is simply partial application.

Rung 3: To manage counterfactuals, we first note that we will consider \emph{$n$-world twinning}, which creates $n$ copies of the original model, because general counterfactual reasoning may require more than 2 copies. We now ask: what is the type of a model under twinning? The twinned version of the car model where only \texttt{fuel} is intervenable and \texttt{carStarts} is returned has type $\IntvMonad{[(\Maybe{\booleans})^n]}{\booleans^n}$: an intervenable computation with $n$ intervention labels (one \texttt{fuel} intervention label per world), and returning $n$ copies of \texttt{carStarts}. The trick is to make each causal model carry its behaviour under twinning, for any world count $n$. In other words the \emph{twinnable} car model is a world-polymorphic vectorised computation of type $\forall n. \IntvMonad{[(\Maybe{\booleans})^n]}{\booleans^n}$, where $\forall$ denotes System-F-style polymorphism \cite{reynolds74_system_f,Girard1972InterpretationFE}. Endogenous variables like \texttt{carStarts} become $n$-dimensional vectors, where $n$ is the number of worlds, and compose together using the applicative interface \cite{gibbons2016naperian}. Twinning is then a simple type application. The corresponding technical setting is the functor category $[\Bij, \catc]$, where $\catc$ is the semantic model of the underlying PPL, e.g.~$\catc = \Qbs$ if it is continuous and higher order. Each type $A$ is informally an $n$-indexed family of $\catc$-objects $\sem{A}(n)$, and each function $A \to B$ is interpreted as an $n$-indexed family of morphisms $\sem{A}(n) \to \sem{B}(n)$, i.e.~its behaviour under $n$-world twinning. This use is similar to names in nominal sets \cite{pitts2013nominalsets}.

We derive $\LHO$ by denotational and functional programming techniques. In \cref{section:densem}, we develop a denotational semantics for our hierarchy of languages, which we show equivalent to the operational description (\cref{thm:L2adequacyfullabs,thm:L3adequacyfullabs}). We then turn the categorical abstractions we develop into functional programming abstractions (\cref{section:higherorder}). 

\subsubsection{Prior existing causal languages and libraries}
Many practical tools like DoWhy (part of the PyWhy ecosystem~\cite{pywhy}) are libraries for manipulating SCMs as graphs and thus have limited expressiveness. DoWhy's GCM module nonetheless reaches rung~3, computing sample-level counterfactuals on fitted mechanisms~\cite{blobaum2024dowhygcm}. Among causal PPLs proper, several stop at rung~2. Brul\'e's Whittemore~\cite{brule2018whittemore}, an early embedded causal DSL, covers rung-2 identification and estimation. Pyro provides a $\opdo$ effect handler~\cite{bingham2018pyro}, the design that ChiRho generalises; PyMC's $\opdo$ performs graph surgery on the computation graph~\cite{abrilpla2023pymc,pymcdo}. Witty et al.~\cite{witty2019synthesis} treat an intervention as a source transformation that replaces subterms, powerful (shift, variance-scaling, mechanism replacement) but non-compositional. In contrast, our languages only expose simple typed holes, which generalise to the same powerful classes of interventions for free at higher-order (\cref{subsection:higherorderintv}). Rung-3 queries can sometimes be re-encoded on such systems by hand by an identifiability argument~\cite{shpitser2012counterfactualstested}.

% Among functional probabilistic languages, 
Four systems support rung 3 natively~\cite{tavares2021language,chirho,perov2020multiverse,kiesel2023whatif}. \textbf{OmegaC}~\cite{tavares2021language} obtains rung 2 through a non-standard lazy evaluation strategy: a let-binding evaluates the bound expression $e$ to a closure $\langle \Gamma, e \rangle$ and $\opdo(x \mapsto e')$ retroactively modifies that closure $\langle \Gamma[x \mapsto e'], e\rangle$, reminiscent of dynamic binding. Rung 3 is enabled by making each program by construction an SCM-like deterministic function of an explicit exogenous noise $\omega$. OmegaC is limited in expressiveness (sampling and building random variables are separate), and only has an operational semantics. We conjecture our formalism can model OmegaC's core denotationally. \textbf{MultiVerse}~\cite{perov2020multiverse} realises Pearl's abduction--action--prediction recipe as an importance-sampling engine, without a formal semantics. \textbf{ChiRho}~\cite{chirho,tutorialchirho}, built from effect handlers and tensor broadcasting, is both the system closest to ours, and our initial motivation. Its recent distillation into typed functional abstractions~\cite{lafi-wang} bears clear similarities to $\LHO$: intervention \emph{points} tracked in a reader monad (our \emph{labels} leave open how an intervention is consumed), vectorisation used in an applicative way for twinning. Two differences are key. First, our rungs 2 and 3 are decoupled: interventions never create worlds; twinning alone does. Under ChiRho's multi-world handler, by contrast, each intervention doubles the number of worlds. Second, $\LHO$ fixes the number of worlds at twinning time and whether a variable is vectorised ($X$ or $\Endo(X)$) determines whether it is exogenous or endogenous. In contrast, ChiRho vectorises its values by default ($\MVal$ in the typed distillation~\cite{lafi-wang}) and encodes an implicit sharing inference using tensor broadcasting: a random draw is shared wherever broadcasting permits, while sample sites downstream of an intervention will draw fresh noise per world. ChiRho's $\MVal$ thus allows finer control over sharing than $\LHO$, but its relation to Pearl's twinning semantics is less clear. Finally, \textbf{probabilistic logic programming} also reaches rung 3: Kiesel et al.~\cite{kiesel2023whatif} answer counterfactual queries on ProbLog programs by a transformation proved correct against Pearl's semantics.

% \subsection{Summary of my contributions}

% ChiRho (LAFI) and semantics (POPL). 

\subsection{Brief overview of notions of causality}

\subsubsection{Logical, Diagrammatic, Categorical notions of Causality}

\paragraph{Functorial Semantics}

\paragraph{Logical characterisation of Ibeling and Icard}

\paragraph{Modulating Normalisation}

\subsubsection{"Causality" in other fields}

\paragraph{Concurrency}

\paragraph{Quantum Computing}

% \section{From quantum protocols to causality}

% \subsection{Foundations of Quantum Programming}
% \begin{itemize}
%     \item QPL, primitives, and semantics in CPTP. 
%     \item Axiomatisations: CQM -- ZX etc; or computing, -- Sam, Chris etc.
% \end{itemize}

% \subsection{Distributed Quantum Processes}
% [Motivation]
% \begin{itemize}
%     \item Quantum protocols (teleportation, BB84, etc.) have multiple quantum computers talking to each other. 
%     \item Distributed quantum computing seen as way of scaling quantum computing. 
%     \item Programmatically we have programs literally communicating to each other. Teleportation for example is a program like [this], which doesn't capture the communication aspect.
% \end{itemize}
% [Approaches and difficulties]
% \begin{itemize}
%     \item To capture these things people have developed several approaches, including quantum process calculi. Grammar looks like [get this from my proposal], semantics given by bisimulaton. But operational, lacks denotational foundation, easily makes something that doesn't necessarily respect physics [give lqCCS' example]. 
%     \item Efforts in making it work but no consensus yet. 
% \end{itemize}
% My master's thesis proposed a different approach. 
% \begin{itemize}
%     \item Algebraic Effects approach: add effects and study interactions in an algebraic theory. Sam's theory = baseline, I added I/O operations.
%     \item Study models, in particular give models that are sensible denotationally as well. 
%     \item Proved these in my masters
%     \item Recently proved completeness as well w help of Claude.  
% \end{itemize}

% \subsection{Causality as the underlying phenomenon}
% [Possible avenues, where I admit I'm not as familiar]

% \subsubsection{Causality and Higher-order quantum computing}
% \begin{itemize}
%     \item Higher order quantum processes = combs, nicely written as these string diagrams with holes. Definition of combs based on non-signalling constraints. 
%     \item Programs with I/O can sort of be thought of as a class of HO quantum process with definite causal order determined by the I/O trace. Indeed, I conjecture this is the case. This seems to suggest that studying quantum comm, and maybe quantum concurrency, will benefit from studying idea of causality.
%     \item Indefinite causal order
% \end{itemize}

% \subsubsection{Causality and Concurrency}
% \begin{itemize}
%     \item "Causal" Models of concurrency, like pomsets. Recently new work on axiomatisation/equational theory like Jack/Cristina's work! 
%     \item Causal models of concurrency with inconsistencies, like event structures.
% \end{itemize}

% % \begin{itemize}
% %     \item Quantum Protocols lit review, esp about ppl not understanding DQC.  
% %     \item Distributed computing, happens-before, causal ordering with Pomsets
% %     \item Higher order quantum computing, briefly, and how they could induce interesting things both on the pomset side (logics).
% % \end{itemize}

\section{Project Proposal}
\subsection{Summary of the current work done}

\subsection{Developing the field of Causal PPLs}
\subsubsection{Putting the language into practice}
\subsubsection{Expressivity}

\subsection{Better understanding our categorical constructions}
\subsubsection{Relation to other constructions}
\subsubsection{A more synthetic approach?}

\subsection{Understanding causality using a graph-centered approach}

\subsection{Understanding relationship with `causality' as used in other fields}
\subsubsection{Concurrency}
\subsubsection{Higher-order processes and quantum causality}


% \subsection{Towards practical causal probabilistic programming}
% \begin{itemize}
%     \item [low risk] Building understanding, e.g. Omega
%     \item [low risk] Putting into practice.
%     \item [med risk] How to implement Cyclic causality and feedback loops, and how they relate to loops/recursion/laziness in programming. 
% \end{itemize}
% \subsection{Understanding DC and QC from the lense of causality.}
% \begin{itemize}
%     \item [low risk] Understanding quantum comm using HOQC
%     \item [med risk] Understanding DQC using process matrices? Clear analogies to be made, understand it categorically. E.g. can I use concurrency things to model quantum switches? Or indeed can I use process matrices generally to do more general concurrency? [Cite Kim as well]
%     \item [higher risk] Categorically understand how DC and HOQC relate in terms of causality? Can I use DC's abstractions to understand QC? Already work in this direction. 
% \end{itemize}
% \subsection{Building a Unified Understanding of Causality?}
% Two extensional ways of characterising causality, with intervention, and with signalling constraints. Both are shown to have non trivial relationship with having an explicit causal structure. How do they relate? I don't understand this and want to understand this better.

\printbibliography
\end{document}
